\documentclass[11pt,a4paper]{article}
\usepackage[T1]{fontenc}
\usepackage[utf8]{inputenc}
\usepackage{lmodern}
\usepackage[margin=28mm]{geometry}
\usepackage{amsmath,amssymb,amsthm,mathtools}
\usepackage{booktabs,array}
\usepackage{microtype}
\usepackage[authoryear,round]{natbib}
\usepackage[colorlinks=true,linkcolor=blue,citecolor=blue,urlcolor=blue]{hyperref}
\hypersetup{pdftitle={Real inflection points of perturbations of concurrent lines},pdfauthor={Oleksiy Babanskyy}}
\newtheorem{theorem}{Theorem}[section]
\newtheorem{proposition}[theorem]{Proposition}
\newtheorem{lemma}[theorem]{Lemma}
\newtheorem{corollary}[theorem]{Corollary}
\theoremstyle{remark}
\newtheorem{remark}[theorem]{Remark}
\newcommand{\R}{\mathbb R}
\newcommand{\C}{\mathbb C}
\newcommand{\RP}{\mathbb {RP}}
\newcommand{\CP}{\mathbb {CP}}
\newcommand{\Hess}{\operatorname{Hess}}
\newcommand{\Infl}{\operatorname{Infl}}
\newcommand{\im}{\operatorname{Im}}
\DeclareMathOperator{\ord}{ord}
\title{Real inflection points of perturbations\\of concurrent lines}
\author{Oleksiy Babanskyy}
\date{}

\begin{document}
\maketitle

\begin{abstract}
Let \(P\) be a real binary form of degree \(n\geq3\) with \(n\)
distinct real projective roots, and let \(V\) be a real homogeneous
form of degree \(n\), nonzero at the concurrent point.
We give an exhaustive criterion for the real inflection points of
\(P+\mu V=0\). If every real projective zero of the second polar
\(V_{zz}\) on each limiting line is simple, then every sufficiently
small nonzero real parameter gives a complex nonsingular curve whose
ordinary real inflection points are in bijection with those zeros.
A uniform estimate excludes real inflection points approaching the
concurrent point. For every fixed \(P\), explicit choices of the polar
realize all even totals from \(0\) to \(n(n-2)\) in even degree and
all totals \(n,n+2,\ldots,n(n-2)\) in odd degree, within one
rigid-isotopy class. In degree four we characterize all fourteen
attainable cyclically labelled line-count vectors, and in even degree
we determine the count on each oval. Radial profiles and the harmonic
Stern family are applications. Two explicit quartics with the same
second polar have six and eight ordinary real flexes, illustrating
the role of simplicity. The complex specialization cycle is identified
as an application of known divisor-specialization theory; the focus
is the exhaustive real criterion and its explicit constructive use.
\end{abstract}

\noindent\textbf{Keywords:} Real algebraic curves; inflection points;
Hessian; deformations; rigid isotopy.

\section{Introduction and main results}

A real inflection point, or flex, of a nonsingular real projective
plane curve is a real point at which the tangent has intersection
multiplicity at least three. The flex is \emph{ordinary} when this
multiplicity is exactly three. We count distinct real projective
points. Nonsingularity below means nonsingularity over \(\C\).

We study a concrete design problem: how can one determine and vary
the real flex count of a perturbation of concurrent lines while
keeping its topology under control? The maximal number \(n(n-2)\)
allowed by Klein's bound and constructions attaining it are
classical; see \citep[Theorem~1.1]{brugallelopez}. Moreover,
Brugall\'e--L\'opez de Medrano already construct odd-degree
maximally inflected curves with all their real flexes on a
pseudoline \citep[Proposition~7.2]{brugallelopez}. Our question
concerns a fixed line configuration and a specified perturbation,
including counts below the maximum.

Fix \(n\geq3\), coordinates \([x:y:z]\) on \(\CP^2\), and
\(O=[0:0:1]\). Let
\begin{equation}\label{eq:family}
 P=\prod_{i=1}^n\ell_i\in\R[x,y]_n,\qquad
 F_\mu=P+\mu V,\qquad C_\mu=\{F_\mu=0\},
\end{equation}
where the real linear factors \(\ell_i\) are pairwise nonproportional
and \(V\in\R[x,y,z]_n\). Put
\begin{equation}\label{eq:polar}
 c=V(O),\qquad S=V_{zz},\qquad L_i=\{\ell_i=0\}.
\end{equation}
The derivative \(V_{zz}\) is the second polar with respect to \(O\)
in these coordinates, up to a nonzero conventional constant.
It is homogeneous of degree \(n-2\).

\begin{theorem}[Real polar criterion]\label{thm:polar}
Assume \(c\ne0\) and that every real projective zero of \(S|_{L_i}\)
is simple, for each \(i\). Write
\[
 r_i=\#\{a\in L_i(\R):S(a)=0\}.
\]
For this fixed pair \(P,V\), there is \(\varepsilon(P,V)>0\) such
that, for every real \(\mu\) with \(0<|\mu|<\varepsilon(P,V)\),
the curve \(C_\mu\) is nonsingular and has exactly
\begin{equation}\label{eq:realcount}
                         N_{\R}(C_\mu)=\sum_{i=1}^n r_i
\end{equation}
real flexes, all ordinary. They are the unique continuations of the
real points of \(P=S=0\). No real flexes approach \(O\).
\end{theorem}

The simplicity assumption concerns only real zeros; multiple nonreal
zeros are permitted. Since \(S(O)=n(n-1)c\ne0\), none of the
restrictions vanishes identically and none has a zero at \(O\).
The statement includes points tending to \(z=0\), without requiring
the flexes themselves to remain at infinity.

\begin{theorem}[Explicit attainable totals]\label{thm:construction}
For every fixed \(P\) as above, the set of totals attainable under
the hypotheses of Theorem~\ref{thm:polar} is
\begin{equation}\label{eq:spectrum}
 \begin{cases}
 \{0,2,4,\ldots,n(n-2)\},& n\ \text{even},\\
 \{n,n+2,n+4,\ldots,n(n-2)\},& n\ \text{odd}.
 \end{cases}
\end{equation}
Explicit perturbations realizing all these totals can be chosen with
\(V(O)=1/[n(n-1)]\). For a sufficiently small common positive
parameter, the resulting curves belong to one rigid-isotopy class.
In odd degree they have one pseudoline and no ovals; in even degree
they have \(n/2\) nonnested ovals.
\end{theorem}

Rigid isotopy here is a path of real degree-\(n\) projective forms
whose complex zero curves are nonsingular. Flexes may become
nonordinary along such a path. Theorem~\ref{thm:construction}
describes the transverse small-perturbation regime of
\eqref{eq:family}; in particular its odd-degree lower endpoint
does not assert a minimum for the entire rigid-isotopy class.
The construction realizes all totals, not arbitrary independent
choices of the vector \((r_1,\ldots,r_n)\). Proposition~\ref{prop:quarticvectors}
classifies these vectors in degree four, and Corollary~\ref{cor:ovalcount}
assigns the flexes to the individual ovals in even degree.

The global existence of all admissible real flex totals is classical:
Ronga's Theorem~II gives every even total from zero and every odd
total from three up to the Klein bound \citep[Theorem~II, p.~197]{ronga}.
Our construction imposes additional constraints: the concurrent line
configuration is prescribed, the perturbations are given explicitly,
and the resulting curves can be chosen in one rigid-isotopy class.
Thus the odd lower endpoint in our family describes its transverse
small-parameter regime, not the global spectrum.

Changes of two real flexes through a contact of order four in generic
families of nonsingular real curves are classical; see
\citep[Section~3.1]{arroyoetal}. Our purpose is to determine the entire
small-parameter real flex locus for a specified perturbation of a
fixed concurrent line configuration, and to construct explicit
profiles realizing the totals in \eqref{eq:spectrum}.

Several ingredients have explicit antecedents. The strict sign of
the binary Hessian is classical \citep{causare}. Second derivatives
of the deformation occur in the determinant in
\citep[Example~5.3]{esteves}, although that example does not directly
apply to our concurrent configuration. Here the leading residual
Hessian is computed directly as \(B_PV_{zz}\). The complex cycle in
Appendix~\ref{app:cycle} follows from
\citep[Theorem~4.1]{esteves} or \citep[Proposition~5.2]{ems}.
The implicit function theorem, coefficient integration and compactness
complete the argument once the uniform real exclusion estimate is proved.

Local theories of inflections under deformation provide complementary
perspectives. Garay develops equivalence and versality with respect
to inflections \citep[Sections~3.2--3.6]{garay}.
Bruce, Fernandes and Tari study local inflection invariants; their
finiteness criterion excludes germs containing a line
\citep[Theorem~3.3(7) in the arXiv version]{bruceetal}.
Our limiting curve is a union of straight lines. We divide the
family Hessian by the parameter and use its residual equation,
together with a uniform real estimate at the concurrent point.
The contribution is the exhaustive real criterion for fixed \(P,V\)
and its explicit constructive use.

\section{The real polar criterion}

\subsection{A nonsingular rescaled model}

\begin{lemma}\label{lem:smooth}
If \(P\) has distinct projective roots and \(V(O)=c\ne0\), then
\(C_\mu\) is nonsingular for every sufficiently small nonzero
complex \(\mu\). No polar-simplicity assumption is needed.
\end{lemma}

\begin{proof}
Write
\[
 V=\sum_{j=0}^n V_j(x,y)z^{n-j},\qquad V_0=c,
\]
with \(V_j\) binary homogeneous of degree \(j\).
For \(\mu=u^n\ne0\), make the projective coordinate change
\([x:y:z]=[uX:uY:Z]\) and divide the equation by \(u^n\).
The transformed form is
\[
 G_u=P(X,Y)+\sum_{j=0}^n u^jV_j(X,Y)Z^{n-j},
 \qquad G_0=P(X,Y)+cZ^n.
\]
A singular point of \(G_0=0\) would have \(Z=0\) and
\(P_X=P_Y=0\). Euler's identity would then give a multiple
projective root of \(P\), a contradiction. The nonsingular locus
in the coefficient space is open, so \(G_u=0\) stays nonsingular
for \(|u|\) small. Every small nonzero complex \(\mu\) has such
an \(n\)-th root.

For real \(\mu=\sigma u^n\), with \(u>0\) and
\(\sigma\in\{1,-1\}\), the same real coordinate change gives instead
\begin{equation}\label{eq:realmodel}
 G_{\sigma,u}
 =P(X,Y)+\sigma\sum_{j=0}^n u^jV_j(X,Y)Z^{n-j},
 \qquad G_{\sigma,0}=P(X,Y)+\sigma cZ^n.
\end{equation}
Both real signs therefore have nonsingular real limiting models.
\end{proof}

The model \(P+cZ^n=0\) also appears in the complex inflection count of
\citet[proof of Theorem~6.8, author-hosted version]{breidingkohnsturmfels}.
For a prescribed perturbation \(V\), the exhaustive real count below
additionally uses the uniform exclusion estimate at the concurrent point.

\subsection{The Hessian calculation}

Set \(B_P=P_{xx}P_{yy}-P_{xy}^2\). The ternary Hessian
\(\mathcal H_\mu=\det\Hess_{x,y,z}F_\mu\) satisfies the exact identity
\begin{equation}\label{eq:hessian}
 \mathcal H_\mu
   =\mu B_PS+\mu^2M+\mu^3\det\Hess V,
\end{equation}
where
\begin{equation}\label{eq:M}
\begin{split}
 M={}&(P_{xx}V_{yy}+P_{yy}V_{xx}-2P_{xy}V_{xy})V_{zz}\\
     &-P_{xx}V_{yz}^2-P_{yy}V_{xz}^2+2P_{xy}V_{xz}V_{yz}.
\end{split}
\end{equation}
Indeed, the \(z\)-row and \(z\)-column of \(\Hess P\) vanish;
expanding the determinant by powers of \(\mu\) gives
\eqref{eq:hessian}--\eqref{eq:M}. Thus
\begin{equation}\label{eq:normalized}
 K_\mu=\mathcal H_\mu/\mu
       =B_PS+\mu M+\mu^2\det\Hess V
\end{equation}
extends polynomially to \(\mu=0\), with \(K_0=B_PS\).
For \(\mu\ne0\), the zeros of \(K_\mu\) on \(C_\mu\) are its flexes.

For completeness, the local relation to contact multiplicity is
especially simple. In an affine chart, let the curve be \(y=h(x)\)
with \(F_y\ne0\), where \(F\) denotes the dehomogenized degree-\(n\)
form. Euler's identities and implicit differentiation give, on the curve,
\begin{equation}\label{eq:contact}
 \det\Hess(F^{\mathrm{hom}})
        =(n-1)^2F_y^3h''.
\end{equation}
Here the Hessian is dehomogenized in the same chart.
The factor multiplying \(h''\) is a unit. A tangent contact of order
\(k\) therefore contributes Hessian intersection multiplicity
\(k-2\), and a transverse Hessian intersection is an ordinary flex.

\subsection{Uniform exclusion at the concurrent point}

We include a proof of the classical binary-Hessian sign to keep the
normalization and strictness explicit; compare
\citep[Proposition~1 and proof of Theorem~1]{causare}.

\begin{lemma}\label{lem:binary}
For the completely real squarefree binary form \(P\),
\(B_P(v)<0\) for every \(v\in\R^2\setminus\{0\}\).
In particular, \(B_P\leq-\eta\) on the unit circle for some \(\eta>0\).
\end{lemma}

\begin{proof}
Choose real linear coordinates putting the point in the chart \(x=1\)
and all roots in that chart, and write
\(f(t)=P(1,t)=a\prod_i(t-a_i)\) with distinct real \(a_i\).
Away from the roots,
\begin{equation}\label{eq:binary}
\begin{split}
 B_P(1,t)&=(n-1)\bigl(nff''-(n-1)(f')^2\bigr)\\
 & =-(n-1)f(t)^2\sum_{i<j}
       \left(\frac1{t-a_i}-\frac1{t-a_j}\right)^2<0.
\end{split}
\end{equation}
At a root the first expression equals
\(-(n-1)^2f'(t)^2<0\). A real invertible linear coordinate change
multiplies the Hessian determinant by the positive square of its
determinant. Rescaling a nonzero vector also preserves the sign,
since \(B_P\) has even degree \(2n-4\). This proves the assertion
in every direction; compactness gives \(\eta\).
\end{proof}

Work near \(O\) in \(z=1\), and write \((x,y)=\rho v\) with
\(\rho>0\) and \(v\in S^1\). Since \(V(\rho v,1)\) tends uniformly
to \(c\ne0\), points of the real curve near \(O\) satisfy
\begin{equation}\label{eq:muscale}
 |\mu|=\left|\frac{P(\rho v)}{V(\rho v,1)}\right|
                 \leq C\rho^n.
\end{equation}
Uniformly in \(v\), we have
\[
 B_P(\rho v)=\rho^{2n-4}B_P(v),\qquad
 S(\rho v,1)=n(n-1)c+O(\rho),\qquad
 M(\rho v,1)=O(\rho^{n-2}).
\]
The second derivatives of \(V\) are bounded in this neighborhood.
Combining these estimates with \eqref{eq:normalized} and
\eqref{eq:muscale}, \emph{on the curve}, gives
\begin{equation}\label{eq:central}
 \frac{K_\mu(\rho v,1)}{\rho^{2n-4}}
 =n(n-1)cB_P(v)+O(\rho)+O(\rho^2)+O(\rho^4).
\end{equation}
The leading term is uniformly separated from zero by
Lemma~\ref{lem:binary}. A fixed sufficiently small real neighborhood
of \(O\) therefore contains no flexes of any small nonzero-parameter
curve. Also \(O\notin C_\mu\), since \(F_\mu(O)=\mu c\).
The estimate is uniform over all approach directions and uses no
assumption on the relative rates at which \(\mu\) and \(\rho\) vanish.

\subsection{Continuation and exhaustion}

Let \(a\) be a real point of \(P=S=0\). It is not \(O\) and belongs
to precisely one \(L_i\). The differential of \(P\) is nonzero
there. Moreover \(B_P(a)\ne0\), and the real simplicity of
\(S|_{L_i}\) makes \(P=K_0=0\) a transverse intersection.
The real analytic implicit function theorem applied in a projective
chart to \(F_\mu=K_\mu=0\) gives exactly one nearby real solution
for every sufficiently small \(\mu\). Its transverse intersection
persists; \eqref{eq:contact} makes the flex ordinary.
This also treats points on \(z=0\).

To exclude other real flexes, suppose that such points occur for a
sequence of nonzero real parameters tending to zero. Compactness
of \(\RP^2\) gives a convergent subsequence. Its limit lies on
\(P=0\), and \eqref{eq:central} excludes \(O\). At any other real
point, \(K_0=B_PS\) and \(B_P\ne0\), so the limit must be a real
zero of \(S|_{L_i}\). The local uniqueness just proved is a
contradiction. Together with Lemma~\ref{lem:smooth}, this proves
Theorem~\ref{thm:polar}.

\section{Constructing all attainable totals}\label{sec:construction}

Write \(d=n-2=2a+\epsilon\) with \(\epsilon\in\{0,1\}\).
For each \(i\), the restriction \(S|_{L_i}\) has projective degree
\(d\). Nonreal roots come in conjugate pairs, counted with
multiplicity, whereas its real roots are simple. Consequently
\[
 r_i\equiv d\pmod2,\qquad \epsilon\leq r_i\leq d.
\]
Theorem~\ref{thm:polar} gives the containment in
\eqref{eq:spectrum}. We now realize every total in that set for
the specified \(P\).

Choose unit vectors \(v_i\) on the \(n\) distinct line directions.
There is a real linear form \(\lambda\) such that
\[
 \lambda_i=\lambda(v_i)^2>0
 \quad\hbox{are pairwise distinct}.
\]
Indeed, it suffices to avoid the finitely many proper linear
conditions
\(\lambda(v_i)=0\), \(\lambda(v_i-v_j)=0\), and
\(\lambda(v_i+v_j)=0\) in the two-dimensional dual space.
The latter vectors are nonzero because the directions are distinct.

Choose pairwise distinct real numbers \(t_1,\ldots,t_a\), none equal
to a \(\lambda_i\), and define
\begin{equation}\label{eq:constructS}
 S=z^\epsilon\prod_{j=1}^{a}
       \bigl(z^2+t_j(x^2+y^2)-\lambda(x,y)^2\bigr).
\end{equation}
On \(L_i\), with coordinates \([x:y:z]=[\rho v_i:z]\), each
factor is \(z^2+(t_j-\lambda_i)\rho^2\). It gives two simple real
roots when \(t_j<\lambda_i\), and none when \(t_j>\lambda_i\).
The factor \(z^\epsilon\) adds \(\epsilon\) simple roots.
Roots belonging to different factors cannot coincide because the
\(t_j\) are distinct. Also \(S(O)=1\).

For one \(t_j\), the number \(b_j=\#\{i:\lambda_i>t_j\}\) can be
any integer from \(0\) to \(n\). Thus
\begin{equation}\label{eq:sumcounts}
                    \sum_i r_i=\epsilon n+2\sum_{j=1}^{a}b_j.
\end{equation}
Every \(K\in\{0,\ldots,an\}\) is a sum of \(a\) integers in
\(\{0,\ldots,n\}\): write \(K=qn+r\), use \(q\) copies of \(n\),
the nonzero remainder if present, and zeros. When \(q=a\), the
remainder is zero. Repeated values of \(b_j\) are realized by
distinct \(t_j\) in the same open interval. Formula
\eqref{eq:sumcounts} therefore gives every desired total.
For \(n=3\) the product is empty, \(S=z\), and the same argument
gives the sole total \(3\).

There is no extension or interpolation problem for these restrictions:
we have constructed a single homogeneous \(S\). If
\(S=\sum_{j=0}^{n-2}s_j(x,y)z^{n-2-j}\), set
\begin{equation}\label{eq:integrate}
 V=\sum_{j=0}^{n-2}
       \frac{s_j(x,y)z^{n-j}}{(n-j)(n-j-1)}.
\end{equation}
Then \(V_{zz}=S\) and \(V(O)=1/[n(n-1)]\).
Terms \(A_{n-1}(x,y)z+B_n(x,y)\) may be added without changing
the limiting criterion. The topology and common-parameter assertions
of Theorem~\ref{thm:construction} follow from the next section.

\begin{table}[ht]
\centering
\begin{tabular}{ccc}
\toprule
Degree & Radial totals & Totals from \eqref{eq:constructS}\\
\midrule
4 & \(0,8\) & \(0,2,4,6,8\)\\
5 & \(5,15\) & \(5,7,9,11,13,15\)\\
6 & \(0,12,24\) & \(0,2,4,\ldots,24\)\\
\bottomrule
\end{tabular}
\caption{Attainable totals in the respective transverse
small-parameter regimes for each fixed line configuration.
The radial column follows from Corollary~\ref{cor:radial}.}
\label{tab:counts}
\end{table}

\begin{proposition}[Quartic line-count vectors]\label{prop:quarticvectors}
Fix four distinct concurrent real lines, labelled in cyclic order
in \(\RP^1\). Among perturbations satisfying
Theorem~\ref{thm:polar}, the attainable vectors
\((r_1,r_2,r_3,r_4)\) are exactly the fourteen members of
\(\{0,2\}^4\) other than
\[
                         (2,0,2,0),\qquad(0,2,0,2).
\]
Each can be realized with \(V(O)=1/12\). The labels are retained:
no quotient by cyclic rotation or permutation is taken.
\end{proposition}

\begin{proof}
Write a general second polar as
\[
 S=az^2+zL(x,y)+Q(x,y),\qquad
 \Delta_S=L^2-4aQ,\qquad a=12V(O)\ne0,
\]
where \(L,Q\) have degrees one and two. For a representative \(v_i\)
of the direction of \(L_i\), the restricted form in coordinates
\([\rho:z]\) is
\(az^2+\rho L(v_i)z+\rho^2Q(v_i)\).
It is nonzero at \(\rho=0\). Simplicity of its real zeros
is equivalent to \(\Delta_S(v_i)\ne0\), and \(r_i=2\)
or \(0\) according as this value is positive or negative.

A nonzero binary quadratic has at most two distinct zeros in
\(\RP^1\). Thus its signs at four cyclic directions cannot
alternate, since alternation would require a zero in each of
four intervening intervals. Conversely, every nonalternating
sign pattern has zero or two cyclic sign changes. A definite
quadratic realizes either constant pattern. For two sign changes,
choose one zero direction in each of the two corresponding open
intervals and take the product of linear forms vanishing there,
with the appropriate overall sign. This realizes the pattern
with nonzero values at all four given directions.

For any such quadratic \(\Delta_S\), set
\[
                S=z^2-\frac{\Delta_S}{4},\qquad
                V=\frac{z^4}{12}-\frac{\Delta_S z^2}{8}.
\]
Then \(V_{zz}=S\), \(V(O)=1/12\), and the discriminant of
the restriction is precisely \(\Delta_S(v_i)\).
Theorem~\ref{thm:polar} proves sufficiency for every fixed \(P\).
\end{proof}

\section{Topology and rigid isotopy}\label{sec:topology}

\begin{proposition}\label{prop:topology}
Fix \(P\), \(c\ne0\), and \(\sigma\in\{1,-1\}\).
For every fixed \(V\) with \(V(O)=c\), the curves
\(P+\mu V=0\), for sufficiently small real \(\mu\) of sign
\(\sigma\), are rigidly isotopic to
\begin{equation}\label{eq:modelcurve}
                         P(X,Y)+\sigma cZ^n=0.
\end{equation}
In odd degree their real locus is one pseudoline; in even degree
\(n=2m\) it consists of \(m\) nonnested ovals.
For a compact segment of such profiles, a common small parameter
gives a path of nonsingular curves. Polar simplicity need not hold
along this path.
\end{proposition}

\begin{proof}
The rescaled path \eqref{eq:realmodel} consists of nonsingular
forms down to \(u=0\), after shrinking the parameter interval.
For \(u>0\), the real coordinate transformation
\([X:Y:Z]\mapsto[uX:uY:Z]\) is connected to the identity through
invertible real transformations. This proves rigid isotopy to
\eqref{eq:modelcurve}.

The model misses \(O\), so projection to \([X:Y]\in\RP^1\)
is defined. For odd \(n\), each direction has exactly one real
value of \(Z\), up to projective scaling. The projection is a
continuous bijection from a compact space to \(\RP^1\), hence a
homeomorphism. A nonsingular odd-degree curve has nontrivial
mod-two projective homology. Its sole real component is therefore
a pseudoline.

For even \(n=2m\), the sign of \(P\) is well-defined on directions
after choosing unit representatives. The function
\(-P/(\sigma c)\) is positive on \(m\) alternate intervals bounded
by consecutive roots in \(\RP^1\). On one such closed interval
\(I\), choose a continuous unit representative \(v(\theta)\) and put
\[
 f(\theta)=\left(-\frac{P(v(\theta))}{\sigma c}\right)^{1/n}.
\]
The branches \(Z=\pm f(\theta)\) meet at the endpoints and bound
\begin{equation}\label{eq:disks}
 D_I=\{[v(\theta):Z]:\theta\in I,\ |Z|\leq f(\theta)\}.
\end{equation}
The strip with its two endpoint fibers collapsed is a closed
disk. Its projective map is injective: distinct interior
directions remain distinct, and a fixed unit representative
identifies no two values of \(Z\). Thus \(D_I\) is an embedded
disk, missing \(O\). The intervals \(I\) are disjoint, including
their endpoints, and so are the disks. Their boundaries give
precisely \(m\) nonnested ovals.

Finally, let \(V_s\), \(s\in[0,1]\), be a segment with
\(V_s(O)=c\). The rescaled forms converge in coefficient space
to \eqref{eq:modelcurve} uniformly in \(s\). Openness of the
nonsingular locus gives a common bound on \(u\).
Hence one small nonzero real \(\mu\) of sign \(\sigma\) makes
the entire segment nonsingular. No simplicity of \(V_{s,zz}\)
is used for this step.
\end{proof}

Choose one profile from Section~\ref{sec:construction} for each
desired total, and join them by finitely many profile segments.
Taking the minimum of their smoothness thresholds and the
endpoint flex thresholds gives a single positive \(\mu\)
realizing all totals and all connecting paths. This completes
Theorem~\ref{thm:construction}. The same finite-choice argument
places representatives of all fourteen vectors in
Proposition~\ref{prop:quarticvectors} in one rigid-isotopy class
at a common sufficiently small positive parameter.

In odd degree the transformation \([X:Y:Z]\mapsto[X:Y:-Z]\)
also connects the two model signs. It belongs to the identity
component of \(\mathrm{PGL}_3(\R)\): multiply its matrix by
\(-1\) and rotate the \(X,Y\)-plane through \(\pi\).
Positive rescaling of \(Z\) adjusts the magnitude of \(c\).
For fixed \(P\), all sufficiently small odd-degree profiles
with \(V(O)\ne0\), of either parameter sign, consequently lie
in the same rigid-isotopy class. Their antipodal lifts to \(S^2\)
are single embedded circles with two complementary domains.
This topological statement does not assert that nonharmonic
profiles are spherical harmonics.

\begin{corollary}[Flexes on each oval]\label{cor:ovalcount}
Assume the hypotheses of Theorem~\ref{thm:polar}, with \(n\)
even, and fix the parameter sign \(\sigma\). For sufficiently
small \(\mu\) of this sign, label each oval by the interval
\(I_{ij}\) between consecutive directions \(i,j\) on which
\(-P/(\sigma c)>0\). The oval has exactly
\[
                              N_{ij}=r_i+r_j
\]
real flexes. In particular \(N_{ij}\) is even and at most
\(2(n-2)\). If the total is \(n(n-2)\), every oval has
\(2(n-2)\) real flexes.
\end{corollary}

\begin{proof}
Write \(\mu=\sigma u^n\), \(u>0\). The rescaled curves
\eqref{eq:realmodel} converge smoothly to
\eqref{eq:modelcurve}, whose ovals have disjoint tubular
neighborhoods. The implicit function theorem and compactness
identify, for all small \(u\), the perturbed real locus
in each neighborhood with a unique normal graph over its
model oval, with no additional real components.

Let a continued flex tend to \(a\in L_i(\R)\setminus\{O\}\).
Choose bounded local homogeneous representatives
\([x(\mu):y(\mu):z(\mu)]\) with a nonzero limiting
\((x,y)\)-pair in direction \(v_i\). In the rescaled coordinates,
\[
 [x(\sigma u^n):y(\sigma u^n):u z(\sigma u^n)]
                      \longrightarrow [v_i:0].
\]
This endpoint belongs to precisely one model oval, namely that
of the unique admissible interval adjacent to direction \(i\).
Consequently all \(r_i\) continued flexes lie on that oval,
including those whose unscaled limit is at infinity.
Its other endpoint contributes \(r_j\). Exhaustion in
Theorem~\ref{thm:polar} gives the formula. Finally,
\(0\leq r_i\leq n-2\) and \(r_i\) is even; at the maximal
total every \(r_i\) equals \(n-2\).
\end{proof}

For \(n=4\), maximal curves in this transverse family therefore have
inflection type \((4,4)\). The global classification of maximally
inflected quartics also includes \((6,2)\) and \((8,0)\)
\citep[Proposition~7.1]{brugallelopez}; the marked line vectors in
Proposition~\ref{prop:quarticvectors} classify a different object.

Reversing the parameter sign changes the pairing of consecutive
directions. It preserves the total but can change the counts
assigned to the resulting ovals. No constancy of these counts
along arbitrary rigid isotopies is asserted.

\section{Radial profiles and the Stern family}\label{sec:radial}

\subsection{The radial criterion in both parities}

Let \(m=\lfloor n/2\rfloor\), and write
\begin{equation}\label{eq:radial}
 Q(q)=\sum_{k=0}^m c_kq^k,\qquad
 V_Q=\sum_{k=0}^m c_k(x^2+y^2)^kz^{n-2k}.
\end{equation}
Its second polar is encoded in the affine chart \(z=1\) by
\begin{equation}\label{eq:D}
\begin{split}
 D_n(q)&=4q^2Q''+(6-4n)qQ'+n(n-1)Q\\
       &=\sum_{k=0}^{m-1}(n-2k)(n-2k-1)c_kq^k.
\end{split}
\end{equation}
In particular \(S(x,y,1)=D_n(x^2+y^2)\).
The coefficient \(c_m\) disappears in both parities.

\begin{corollary}\label{cor:radial}
Assume \(c_0c_{m-1}\ne0\) and every positive root of \(D_n\)
is simple. If \(r\) is the number of positive roots, then for
every sufficiently small nonzero real \(\mu\), the curve
\(P+\mu V_Q=0\) is nonsingular and has exactly
\[
 \begin{cases}
 n(2r+1),&n=2m+1,\\
 2nr,&n=2m
 \end{cases}
\]
ordinary real flexes. The finite limiting points have squared
radii equal to the positive roots of \(D_n\). In odd degree
there is also one flex tending to infinity on each line; in
even degree there are no such limiting flexes.
\end{corollary}

\begin{proof}
On each line, the affine polar has two simple roots for each
positive root of \(D_n\). In odd degree its homogenization has
a simple factor \(z\) at infinity, with nonzero residual
coefficient \(6c_{m-1}\). In even degree the value at a real
infinite direction is \(2c_{m-1}(x^2+y^2)^{m-1}\ne0\).
Theorem~\ref{thm:polar} applies.
\end{proof}

This gives all radial totals in Table~\ref{tab:counts}. More
generally, given \(D_n=\sum_{k=0}^{m-1}d_kq^k\) with
\(d_0d_{m-1}\ne0\), choose
\begin{equation}\label{eq:radialinverse}
 c_k=\frac{d_k}{(n-2k)(n-2k-1)}\quad(k<m),\qquad
 c_m\ \text{arbitrary}.
\end{equation}
Choosing \(r\) distinct positive and \(m-1-r\) distinct negative
roots realizes any \(r=0,\ldots,m-1\), and prescribes the
limiting squared radii. Repeated negative or nonreal roots are
allowed. Common roots of \(Q\) and \(D_n\) are also allowed when
the positive roots of \(D_n\) remain simple.

For odd degree the real locus is a pseudoline by
Proposition~\ref{prop:topology}, now without any equal-spacing
assumption. For even degree, Corollary~\ref{cor:ovalcount}
gives \(4r\) flexes on each oval, since every line contributes
\(2r\). The fixed-\(P\) isotopy statement of
Proposition~\ref{prop:topology} applies to both radial and
nonradial profiles.

\subsection{Equal spacing, exact locations and triple roots}

For odd \(n=2m+1\), put \(W_n=\im(x+iy)^n\).
For the affine equation \(W_n+\mu Q(x^2+y^2)=0\), let
\[
 T(q)=4\{2q(Q')^3+n(n-3)Q(Q')^2-n^2Q^2Q''\},\qquad
 L_\mu(q)=n^2q^{n-2}D_n(q)+\mu^2T(q).
\]
At real curve points away from the origin, the tangent-Hessian
numerator is
\begin{equation}\label{eq:sectoralK}
                          N=\mu L_\mu(q).
\end{equation}
Here \(N=f_y^2f_{xx}-2f_xf_yf_{xy}+f_x^2f_{yy}\), with
\(f=W_n+\mu Q(x^2+y^2)\), is the affine Hessian evaluated
on the tangent vector \((-f_y,f_x)\).
To check the formula, write \(W_n=\rho^np(\theta)\),
use \(p''=-n^2p\), \((p')^2=n^2(1-p^2)\), and substitute
\(\rho^np=-\mu Q\) in the tangent quadratic form of the
affine Hessian. For example, with \(A=2qQ'-nQ\), the cubic
coefficient simplifies by
\[
             A^3+n^2QA^2-n^2Q^2D_n=q^2T.
\]
The curve equation is feasible exactly when
\(\Delta_\mu(q)=q^n-\mu^2Q(q)^2\geq0\).
Near a simple positive root \(\alpha\) of \(D_n\), this
inequality is strict for small \(\mu\). Its \(2n\) flexes
have the common squared radius
\begin{equation}\label{eq:radialdrift}
 q_\alpha(\mu^2)=\alpha-
   \frac{T(\alpha)}{n^2\alpha^{n-2}D_n'(\alpha)}\mu^2+O(\mu^4),
\end{equation}
determined by \(L_\mu=0\). Their angles satisfy
\(\sin(n\theta)=-\mu Q(q_\alpha)/q_\alpha^{n/2}\).
Local uniqueness accounts for all \(2n\) points.

The \(n\) remaining flexes in this odd sectoral family lie
exactly at infinity. After rotating any limiting direction to
\([1:0:0]\), the polynomial \(W_n(1,y)\) is odd in \(y\).
The radial perturbation is even in \(y\) and odd in \(z\).
The implicit graph \(y=h(z,\mu)\) is therefore odd in \(z\),
so \(h''(0,\mu)=0\); ordinary contact follows from
Corollary~\ref{cor:radial}. Section~\ref{sec:asymmetric}
shows why these exact locations do not survive general
line configurations.

There is a simple exact-location observation that does not
require equal spacing. Suppose \(Q\) has a root of multiplicity
exactly three at \(\alpha>0\). At each of the \(2n\) intersections
of \(P=0\) with \(x^2+y^2=\alpha\), \(Q=Q'=Q''=0\).
The curve is regular there because its transverse derivative
equals that of \(P\). Its tangent is the corresponding
limiting line, whose restriction of the equation is
\(\mu Q(x^2)\) in an orthonormal line coordinate. For
\(\mu\ne0\), the contact order is exactly three. These points
are fixed ordinary flexes, rather than hyperflexes. Consistently,
\(D_n(\alpha)=0\) and
\(D_n'(\alpha)=4\alpha^2Q'''(\alpha)\ne0\).

The following degree-seven profiles illustrate the permitted
degeneracies and the distinction between roots of \(Q\) and \(D_7\):
\begingroup
\renewcommand{\arraystretch}{1.35}
\[
\begin{array}{c|c|c}
 Q(q)&D_7(q)&N_{\R}\\ \hline
 \frac1{42}+\frac q{10}+\frac{q^2}{6}&(q+1)^2&7\\
 1-7q^2+q^3&42(1-q^2)&21\\
 (q-1)^3&-6(q-1)(3q-7)&35
\end{array}
\]
\endgroup
The first has vanishing top radial coefficient and a double
negative root of \(D_7\). In the last, \(q=1\) is a fixed
circle of fourteen ordinary flexes for \(W_7\); the other
positive root of \(D_7\) is \(7/3\).

\subsection{Harmonic specialization}

Let \(L_n\) be the Legendre polynomial normalized by \(L_n(1)=1\),
write \(R^2=x^2+y^2+z^2\), and put
\[
 Z_n=R^nL_n(z/R).
\]
This expression is a homogeneous polynomial. For odd \(n=2m+1\),
the classical sectoral and zonal harmonics give the Stern family
\begin{equation}\label{eq:stern}
                         W_n+\mu Z_n=0;
\end{equation}
see B\'erard--Helffer \citep{berardhelffer}.
The lowering identity
\begin{equation}\label{eq:lowering}
                  (Z_n)_{zz}=n(n-1)Z_{n-2}
\end{equation}
can be checked without importing a new geometric assertion.
The generating function
\[
 \sum_{k\geq0}Z_kt^k
  =\bigl(1-2zt+(x^2+y^2+z^2)t^2\bigr)^{-1/2}
\]
satisfies
\(\partial_z\sum Z_kt^k=t\sum Z_kt^k+
t^2\partial_t\sum Z_kt^k\).
Comparing coefficients gives \(\partial_zZ_k=kZ_{k-1}\);
applying it twice proves \eqref{eq:lowering}.

In the chart \(z=1\),
\[
 Q_n(q)=(1+q)^{n/2}L_n((1+q)^{-1/2}),\qquad
 D_n(q)=n(n-1)Q_{n-2}(q).
\]
Equivalently, the explicit radial coefficients are
\[
 Q_n(q)=\sum_{k=0}^{m}
       \frac{(-1)^k n!}{4^k(k!)^2(n-2k)!}\,q^k.
\]
The odd Legendre polynomial \(L_{n-2}\) has \(m-1\) positive
simple roots \(\tau_j\in(0,1)\); its remaining central root is
zero. Hence \(Q_{n-2}\) has exactly \(m-1\) positive simple
roots \(\tau_j^{-2}-1\), and its leading coefficient is
nonzero. Also \(Q_n(0)=1\). Corollary~\ref{cor:radial} gives
exactly \(n(n-2)\) ordinary real flexes for all sufficiently
small nonzero real \(\mu\). The finite limiting squared radii
are \(\tau_j^{-2}-1\).

For clarity, the elementary Legendre root property used here
follows from orthogonality on \([-1,1]\): if a degree-\(k\)
Legendre polynomial had fewer than \(k\) sign-changing roots
in \((-1,1)\), multiplication by the product of their linear
factors would give a polynomial of degree less than \(k\)
with a nonzero integral against it, contradicting orthogonality.
Thus all \(k\) roots are simple and internal; parity supplies
the positive-root count.

All these real flexes lie on the sole pseudoline, and its
spherical lift has two complementary domains. The nodal
construction of the harmonic family is classical
\citep[Proposition~1(i), preprint version]{berardhelffer};
the present specialization concerns the exhaustive flex
calculation. Maximal existence and concentration on a
pseudoline are not asserted as new.

\section{Examples and the limits of the criterion}

\subsection{A nonradial quartic with six real flexes}

Take
\begin{equation}\label{eq:quartic}
 P=y(x-y)(x-2y)(x-3y),\qquad
 V=\frac{z^4}{12}+\frac{(-x^2+2y^2)z^2}{2}.
\end{equation}
Then \(S=z^2-x^2+2y^2\) has respectively \(2,0,2,2\)
real projective zeros on the four lines. Theorem~\ref{thm:polar}
therefore gives six ordinary real flexes at sufficiently
small nonzero parameter. This total is unavailable to
the transverse degree-four radial profiles of
Corollary~\ref{cor:radial}, which give zero or eight.

The fixed value \(\mu=1/1000\) is also certified directly.
Multiplying \(F_\mu\) by \(12000\) gives
\[
 F=12000y(x-y)(x-2y)(x-3y)+z^4-6x^2z^2+12y^2z^2.
\]
Let \(f=F(x,y,1)\) and let \(h\) be the dehomogenized
primitive ternary Hessian. Elimination in \(y\) gives a
squarefree polynomial \(E(x)\) of degree \(24\). The
penultimate subresultant is linear in \(y\), with coefficient
coprime to \(E\); the leading coefficient of \(f\) in \(y\)
is the nonzero constant \(-72000\). Sturm's method gives
exactly six real roots of \(E\), hence six real affine
intersection points. There are \(24\) distinct complex
affine intersections, the full B\'ezout degree, so all are
simple. Separate exact checks show that the curve is
nonsingular and that none of its four points at infinity
is a flex. The accompanying script
\texttt{examples/check\_quartic.py} reconstructs these
polynomials and checks, using rational arithmetic.
This fixed-parameter certificate is distinct from the
small-parameter theorem.

\subsection{An asymmetric quintic and moving locations}\label{sec:asymmetric}

Let
\[
                   P=y(x-y)(x-2y)(x-3y)(x-4y).
\]
Near \(y=0\), write \(P=d x^4y+e x^3y^2+O(y^3)\);
here \(d=1\), \(e=-10\). For \(Q=1-q\), the polynomial
\(D_5=20-6q\) has positive root \(\alpha=10/3\).
The radial corollary gives fifteen flexes. The local graph
\(y=\mu h_1(x)+\mu^2h_2(x)+O(\mu^3)\) has
\[
 h_1=-\frac{1-x^2}{d x^4},\qquad
 h_2=-\frac{e(1-x^2)^2}{d^3x^9}.
\]
At \(x_0=\pm\sqrt{10/3}\),
\[
 h_1'''(x_0)=\frac{12}{d}x_0^{-5},\qquad
 h_2''(x_0)=-\frac{50e}{d^3}x_0^{-11}.
\]
Solving \(h_{xx}=0\) gives
\(x_{\rm fl}=x_0+25e\,x_0^{-6}\mu/(6d^2)+O(\mu^2)\).
Consequently the squared radii of the two flexes near this line are
\begin{equation}\label{eq:split}
 q_+(\mu)=\frac{10}{3}-\frac{3\sqrt{30}}4\mu+O(\mu^2),
 \qquad
 q_-(\mu)=\frac{10}{3}+\frac{3\sqrt{30}}4\mu+O(\mu^2).
\end{equation}
They have a common limiting radius but distinct radii for
small nonzero parameter.

For the same \(P\), take \(Q=1+q+q^2\). Now \(D_5=20+6q\),
so the five real flexes all tend to infinity. In the chart
\(x=1\) at \([1:0:0]\), write \(y=h(z,\mu)\).
Implicit differentiation gives
\[
                  h_{zz}(0,\mu)=20\mu^2,\qquad
                  z_{\rm fl}(\mu)=\frac{10}{3}\mu+O(\mu^2).
\]
This flex therefore leaves \(z=0\).
Both effects follow by substitution and differentiation;
no count is inferred from a numerical plot.

Intermediate quintic totals also have short formulas.
For this \(P\), the polars
\[
                 S=z(z^2-x^2+\lambda y^2)
\]
with \(\lambda=2,5,10,17\) have respectively \(13,11,9,7\)
simple real zeros on the line configuration. Integrating
twice in \(z\) and applying Theorem~\ref{thm:polar} realizes
these totals. These are statements at sufficiently small
parameter, not additional fixed-parameter certificates.

\subsection{The same polar with six or eight real flexes}\label{sec:samepolar}

The polar alone need not determine the real count when a restricted
root is multiple. Consider
\begin{equation}\label{eq:samepolar}
 P=y(x-y)(x-2y)(x-3y),\qquad
 V_b=\frac{z^4}{12}+\frac{(y^2-x^2)z^2}{2}+bx^4.
\end{equation}
For every \(b\), we have \(V_b(O)=1/12\) and
\[
                         V_{b,zz}=z^2-x^2+y^2.
\]
This polar has two simple real zeros on each of
\(y=0,x=2y,x=3y\), and a double zero at \(a=[1:1:0]\)
on \(x=y\). We show that for sufficiently small \(\mu>0\),
the choice \(b=1\) gives six ordinary real flexes and
\(b=-1\) gives eight.

In the chart \(x=1\), put \(p(y)=y(1-y)(1-2y)(1-3y)\).
Since \(p'(1)=-2\), the local curve is \(y=h(z,\mu)\),
where \(h\) is analytic and even in \(z\). Write
\(h=1+\mu h_1+\mu^2h_2+O(\mu^3)\). Substitution gives
\[
 \begin{split}
 h_1&=\frac b2+\frac{z^4}{24},\\
 h_2&=-\frac{9b^2}{8}-\frac{3bz^4}{16}
       +\frac{bz^2}{4}-\frac{z^8}{128}+\frac{z^6}{48}.
 \end{split}
\]
The normalized flex equation is therefore
\begin{equation}\label{eq:doubleunfold}
 \frac{h_{zz}}{\mu}
             =\frac{z^2}{2}+\frac{b\mu}{2}
                       +O(\mu z^2+\mu^2).
\end{equation}
The quotient extends analytically to \(\mu=0\).
By parity it is analytic in \(w=z^2\), with derivative
in \(w\) close to \(1/2\). For \(b=1\) it is positive
for real \(z\) in a fixed small neighborhood and small
\(\mu>0\). For \(b=-1\) it has exactly one nearby
zero in \(w\), namely \(w=\mu+O(\mu^2)>0\), yielding
two real zeros
\(z=\pm\sqrt{\mu}+O(\mu^{3/2})\).
They are simple zeros of \(h_{zz}\), hence ordinary flexes.

The six simple polar zeros continue as ordinary real flexes.
The nonsingularity lemma and the central exclusion estimate
do not require polar simplicity. Together with projective
compactness and the above local calculation, they exclude
all further real flexes for these fixed profiles.

At \(b=0\), \(a\) is a fixed smooth point with tangent
\(y=x\). Restricting the equation to that tangent gives
\(\mu z^4/12\), so the contact order is exactly four.
Locally \(h-1\) is divisible by \(z^4\), and
\(h_{zz}/(\mu z^2)\) extends to a nonzero analytic
function with value \(1/2\) at \((0,0)\).
Thus there are six ordinary real flexes and this one
hyperflex: seven distinct points, with total real Hessian
intersection multiplicity eight.

There are also exact certificates at \(\mu=1/1000\).
For \(b=1,-1\), respectively, the integer forms
\[
 12000P+z^4+6(y^2-x^2)z^2+12bx^4
\]
have six and eight ordinary real flexes. The accompanying
\texttt{examples/check\_extensions.py} verifies complex
nonsingularity, no flex at infinity, a squarefree elimination
polynomial of degree \(24\), coprimality of the linear
subresultant coefficient, and the real-root count by Sturm's
theorem. These conditions identify all \(24\) distinct
complex curve--Hessian intersections and certify ordinarity.
The exact output is \texttt{EXTENSION\_CHECKS.json}.

The segment \(b\in[-1,1]\) has fixed central value, so
Proposition~\ref{prop:topology} joins these profiles in one
rigid-isotopy class at a common sufficiently small parameter.
The two certificates at \(1/1000\) do not certify the whole
segment at that value. The change through an order-four
contact is classical \citep[Section~3.1]{arroyoetal}; here
the example keeps the second polar fixed and exhibits the
effect of a coefficient lost under two differentiations.

\subsection{Colliding directions and tropical comparison}

The example above concerns a multiple restricted polar root.
It is distinct from a triple root of \(Q\) in Section~5.2,
where the relevant root of \(D_n\) may still be simple.

Uniformity holds on the compact profile segments in
Proposition~\ref{prop:topology}. It is not asserted for
unbounded coefficient families or colliding line directions.
If two directions merge, the rescaled model can become singular
and the strict binary-Hessian bound can degenerate. The theorem
therefore supplies no positive common threshold or assertion
about persistence of distinct flexes through that limit.

For the natural radial degeneration \(\mu=t\), after a generic
rotation of \(x,y\), the degree-\(n\) edge coefficients have
valuation zero and the nonzero lower radial coefficients have
valuation one. The lower Newton hull has a coarse triangular
cell of normalized area \(n^2\). It is not the unimodular
subdivision assumed in the nonsingular tropical lifting
statements of \citep[Theorems~5.6--5.7]{brugallelopez}.
Those particular hypotheses do not apply directly to this
presentation. This observation does not rule out other
tropical presentations and is not evidence of novelty.
Our real criterion and construction require no tropical
lifting assertion.

\appendix
\section{The weighted complex specialization cycle}\label{app:cycle}

For \(\mu\ne0\), let \(\Infl(C_\mu)\) be the intersection
cycle of the nonsingular curve with its ternary Hessian,
including multiplicities. With \(S=V_{zz}\) and \(c\ne0\),
\begin{equation}\label{eq:cycle}
 \lim_{\mu\to0}\Infl(C_\mu)
                  =2n(n-2)[O]+(P\cdot S).
\end{equation}
Here \(P\cdot S\) denotes the projective intersection cycle,
of degree \(n(n-2)\), supported away from \(O\). No
ordinary-flex assumption on its nonreal points is made.

We give a direct specialization argument. The forms \(P\)
and \(K_0=B_PS\) have no common component: \(S(O)\ne0\)
excludes any \(L_i\) from \(S\), and the restriction of \(B_P\)
to a simple-root direction of \(P\) is nonzero.
The family
\[
 \mathcal Z=\{F_\mu=K_\mu=0\}\subset\CP^2\times\Delta
\]
over a sufficiently small complex disk has finite special
fiber. Properness and upper semicontinuity of fiber dimension
make the projection finite after restricting \(\Delta\).
At a point above \(\mu=0\), the three elements
\(\mu,F_\mu,K_\mu\) form a system of parameters in the
regular local ambient ring of dimension three. They are
a regular sequence, also in the order \(F_\mu,K_\mu,\mu\).
Thus the finite family has no \(\mu\)-torsion and is flat
over the one-dimensional smooth base. Its special
intersection cycle is the specialization of the generic one.

On each \(L_i\), \(B_P\) is a nonzero constant times
\(\rho^{2n-4}\), where \([\rho:z]\) are line coordinates
and \(O=[0:1]\). It contributes \(2n-4\) at \(O\).
Summing over the \(n\) lines gives \(2n(n-2)\); the
remaining cycle is \(P\cdot S\), proving \eqref{eq:cycle}.
The total degree is \(3n(n-2)\), as expected.
By \eqref{eq:central}, the central cluster is entirely
nonreal for small real nonzero parameter, with multiplicities
understood.

Formula \eqref{eq:cycle} is an application of existing
specialization theory. On the flat family
\(X=\{P+tV=0\}\), put \(H=\det\Hess(P+tV)=tK\).
The divisor identity is
\[
                 \operatorname{div}(H)=X_0+\operatorname{div}(K).
\]
Since \(K_0\) contains no component of \(P=0\), Theorem~4.1
of \citet{esteves} applies with \(E_i=0\), \(p_i=1\) and
\(F_i=\operatorname{div}(K)\). Equivalently, Proposition~5.2
of \citet{ems} applies with \(G=H\), \(L_i=1\), \(M_i=K\).
The extra irreducibility assumption in
\citep[Example~5.3]{esteves}, and the nondegenerate-system
assumptions of the specialized ramification formulas in
\citep{ems}, do not obstruct these general applications.
The complex cycle alone cannot distinguish real flexes
approaching \(O\) from conjugate pairs. That distinction
is supplied by the uniform real estimate in
Section~2.

\section*{AI assistance}
OpenAI Codex contributed substantially to mathematical exploration and
proof development, symbolic calculations and verification code, literature
searches, and manuscript drafting and revision. The author directed the
project and is responsible for the claims. The author reports using the
most capable GPT model available to him in Codex at the time; exact
historical model-version identifiers are not documented.
The accompanying scripts reproduce the finite example checks. Those
checks do not replace the general proofs, and AI review does not
constitute independent specialist review or formal verification.

\paragraph{License.}
The manuscript and its sources are licensed under
\href{https://creativecommons.org/licenses/by/4.0/}{Creative Commons
Attribution 4.0 International (CC BY 4.0)}. Original companion software
and repository documentation are licensed separately under MIT.

\clearpage
\bibliographystyle{plainnat}
\bibliography{references}
\end{document}
